\begin{pastexam}[label={px:ntu-112-ee-ds-09}]{臺大電機所 112 學年度 第 9 題}
（是非；命題正確填 (A)，錯誤填 (B)；答對得 4 分，答錯倒扣 4 分，未填答不計分，倒扣至是非題總分 0 分為止）

A \textit{cube} can be represented as a \textit{planar graph} and also a \textit{bipartite graph}.
\end{pastexam}

\begin{pastsoln}
答案是 (A)，正確。

立方體的圖 $Q_3$：$8$ 個頂點是立方體的頂點，用三位元 $000$ 到 $111$ 表示；兩個頂點相鄰若且唯若它們只差
一個位元（立方體的 $12$ 條稜）。

\textbf{平面圖}：把立方體從正面往裡看，前面那個正方形畫大、後面那個正方形畫小放在中間，四條連接前後的稜
連接兩個正方形對應的角。 這樣畫 $12$ 條邊都不交叉，所以是平面圖。

\textbf{二分圖}：依位元中 $1$ 的個數分成偶數組 $\{000, 011, 101, 110\}$ 和奇數組
$\{001, 010, 100, 111\}$。 相鄰的兩點只差一個位元，$1$ 的個數奇偶性不同，所以每條邊都一端在偶數組、一端
在奇數組，是二分圖（資結書定理~21.9：二分圖若且唯若沒有奇環；$Q_3$ 的環長都是偶數）。

兩個性質同時成立，正確。
\end{pastsoln}
